\documentclass[conference]{IEEEtran}
% --- Paketler ---
\usepackage[utf8]{inputenc}
\usepackage[T1]{fontenc}
\usepackage[turkish]{babel}
\shorthandoff{=}          % babel-turkish '=' işaretini aktif yapar; tablolarda sorun çıkmasın
\usepackage{amsmath,amssymb}
\usepackage{graphicx}
\usepackage{booktabs}
\usepackage{tabularx}
\usepackage{array}
\usepackage{url}
\usepackage[hidelinks]{hyperref}
\usepackage{enumitem}

\newcolumntype{L}{>{\raggedright\arraybackslash}X}
\newcommand{\todo}[1]{\textcolor{red}{[#1]}}
\usepackage{xcolor}

\begin{document}

\title{ALL/AML Lösemi Alt Tip Sınıflandırması:\\ Literatür, Veri Setleri ve Keşifsel Veri Analizi\\ (Ara Rapor)}

\author{\IEEEauthorblockN{Mehmet Ali \todo{Soyad}}
\IEEEauthorblockA{Öğrenci No: \todo{doldurulacak}\\
Makine Öğrenmesi ve Anomali Tespiti\\
Video: \todo{URL}\quad Kod: \todo{GitHub URL}}}

\maketitle

\begin{abstract}
Akut lenfoblastik lösemi (ALL) ve akut miyeloid lösemi (AML) farklı tedavi protokolleri gerektiren iki lösemi alt tipidir; bu nedenle doğru ayrımları klinik açıdan kritiktir. Bu ara rapor, gen ifadesi mikrodizi verisinden ALL/AML alt tip sınıflandırması yapacak bir makine öğrenmesi çalışmasının temelini atmaktadır. Çalışmada Golub ve ark. (1999) verisi (72 örnek, 7.129 prob) ana veri seti, farklı nesil bir Affymetrix platformunda üretilmiş çok merkezli MILE çalışması (GEO GSE13159) ise dış doğrulama veri seti olarak belirlenmiştir. Golub verisini kullanan çalışmaların taranması, öznitelik seçiminin çoğu zaman çapraz doğrulama döngüsü dışında yapıldığını, nested CV kullanılmadığını ve sonuçların büyük ölçüde yalnızca doğruluk ile raporlandığını göstermiştir. Sızıntıyı düzelten çalışmalarda doğruluk \%92--95 düzeyine inerken sızıntılı tasarımlar \%100 bildirmektedir. \todo{EDA bulgularının 1--2 cümlelik özeti.} Finalde sızıntısız nested CV ile en az üç yöntem karşılaştırılacak ve bulguların platformlar arası genellenebilirliği test edilecektir.
\end{abstract}

\begin{IEEEkeywords}
lösemi, ALL/AML sınıflandırması, mikrodizi, gen ifadesi, öznitelik seçimi, veri sızıntısı, nested çapraz doğrulama
\end{IEEEkeywords}

% =====================================================================
\section{Giriş}

\subsection{Problem ve Klinik Önem}
Akut lösemiler, kökenlerine göre lenfoid öncül hücrelerden gelişen ALL ve miyeloid öncül hücrelerden gelişen AML olarak iki ana gruba ayrılır~\cite{golub1999}. Bu ayrım klinik açıdan kritiktir: iki alt tip farklı kemoterapi rejimleriyle tedavi edilir ve yanlış sınıflandırma hem tedavi başarısını düşürür hem de gereksiz toksisiteye yol açar~\cite{golub1999}. Rutin tanı; morfoloji, immünofenotipleme ve sitogenetik gibi uzmanlık gerektiren birden fazla testin birlikte yorumlanmasına dayanır. Gen ifadesi profillemesi, bu ayrımı tek bir moleküler testle ve nesnel ölçütlerle yapma imkânı sunar~\cite{golub1999,haferlach2010}.

Bu çalışmada problem, \textbf{ALL ve AML alt tip sınıflandırması} olarak tanımlanmıştır. Ders modülündeki ``cancer / non-cancer'' ifadesi biyolojik olarak yanlıştır; her iki sınıf da lösemi alt tipidir.

\subsection{Mikrodizi Teknolojisi}
DNA mikrodizileri, binlerce genin mRNA düzeyini tek bir deneyde ölçen yonga tabanlı bir teknolojidir. Hastadan alınan örnekteki mRNA etiketlenir ve yonga üzerindeki bilinen dizilere sahip problarla hibridize edilir; her probdaki sinyal şiddeti ilgili genin ifade düzeyini yansıtır. Golub verisinde kullanılan Affymetrix Hu6800 yongası 7.129 prob içerir~\cite{golub1999}. Bu teknoloji örnek sayısından çok daha fazla öznitelik üretir ($p \gg n$). Bu durum aşırı öğrenme riskini artırır ve öznitelik seçimini zorunlu kılar~\cite{saeys2007}.

\subsection{Araştırma Soruları}
\begin{itemize}[leftmargin=*]
  \item \textbf{AS1:} Öznitelik seçimi çapraz doğrulama döngüsü içine alındığında (nested CV), Golub verisinde elde edilen performans literatürde sıkça bildirilen \%100 doğruluktan anlamlı ölçüde düşük müdür?
  \item \textbf{AS2:} Golub verisinde eğitilmiş modellerin performansı, farklı nesil bir platformda (HG-U133 Plus 2.0) üretilmiş bağımsız MILE verisine aktarıldığında korunur mu?
  \item \textbf{AS3:} Sınıf dengesizliği (\%65 ALL / \%35 AML) altında doğruluk ile dengeli doğruluk ve ROC-AUC arasında yöntemleri farklı sıralayacak kadar fark var mıdır?
\end{itemize}

\subsection{Katkılar}
\begin{itemize}[leftmargin=*]
  \item Golub verisi üzerindeki çalışmaların doğrulama protokollerini ve sızıntı risklerini karşılaştıran sistematik bir literatür özeti.
  \item Golub ve MILE verilerinin kalite, dağılım ve batch farkları açısından karşılaştırmalı keşifsel analizi (EDA).
  \item Finalde sızıntısız, tekrarlanabilir bir değerlendirme protokolü ve platformlar arası dış doğrulama için hazırlanmış bir kod deposu.
\end{itemize}

% =====================================================================
\section{Literatür Taraması}

\subsection{Arama Protokolü}
Tarama IEEE Xplore, PubMed ve Google Scholar üzerinde 1999--2026 aralığında yapılmıştır; son beş yıl (2021--2026) ayrıca taranmıştır. Kullanılan sorgular ve döndürdükleri sonuç sayıları Tablo~\ref{tab:sorgu}'de verilmiştir.

\begin{table}[htbp]
\caption{Arama sorguları ve sonuç sayıları}
\label{tab:sorgu}
\centering
\footnotesize
\begin{tabularx}{\columnwidth}{@{}c L c c@{}}
\toprule
\# & Sorgu & IEEE & Scholar \\
\midrule
1 & ``Golub'' AND leukemia AND (classification OR ``machine learning'') & \todo{} & \todo{} \\
2 & ``ALL'' AND ``AML'' AND ``gene expression'' AND (classifier OR ``feature selection'') & \todo{} & \todo{} \\
3 & microarray AND ``feature selection'' AND ``small sample'' AND cancer & \todo{} & \todo{} \\
4 & leukemia AND ``gene expression'' AND (``deep learning'' OR transformer) & \todo{} & \todo{} \\
5 & (``data leakage'' OR ``selection bias'') AND microarray AND ``cross-validation'' & \todo{} & \todo{} \\
\bottomrule
\end{tabularx}
\end{table}

\textbf{Dahil etme ölçütleri:} hakemli dergi veya konferans yayını; mikrodizi ya da RNA-seq ile lösemi veya kanser sınıflandırması; doğrulama yöntemi ve sayısal sonuç raporlanmış olması.
\textbf{Dışlama ölçütleri:} tam metnine erişilemeyen, doğrulama yöntemi belirtilmemiş, tekrar yayınlanmış veya DOI'si doğrulanamayan çalışmalar. Sonuçta 22 hakemli yayın ve bir ön baskı seçilmiştir; bunların 11'i son beş yıla aittir.

\subsection{Temel Çalışmalar ve Yöntem Ailesi}
Golub ve ark.~\cite{golub1999}, sinyal-gürültü oranıyla seçilen 50 gen ve ağırlıklı oylama yöntemiyle ALL/AML ayrımını yapmış; ayrıca sınıf bilgisi kullanılmadan kümelemenin bu iki sınıfı kendiliğinden ayırdığını göstermiştir. Sonraki çalışmalar bu veriyi bir kıyaslama standardına dönüştürmüştür. Furey ve ark.~\cite{furey2000} ve Guyon ve ark.~\cite{guyon2002} destek vektör makinelerini (SVM) öne çıkarmış; Guyon'un SVM-RFE yöntemi gömülü öznitelik seçiminin klasiği olmuştur. Dudoit ve ark.~\cite{dudoit2002} basit yöntemlerin (DLDA, k-NN) karmaşık yöntemlerden geri kalmadığını, Díaz-Uriarte ve Alvarez de Andrés~\cite{diaz2006} ise rastgele ormanların (RF) diğer yöntemlerle karşılaştırılabilir performans verdiğini göstermiştir. Saeys ve ark.~\cite{saeys2007} öznitelik seçimini filtre, sarmalayıcı ve gömülü yöntemler olarak sınıflandırmıştır.

\subsection{Doğrulama ve Veri Sızıntısı}
Ambroise ve McLachlan~\cite{ambroise2002}, genler tüm veri üzerinde seçildikten sonra yapılan çapraz doğrulamanın hatayı sistematik olarak düşük gösterdiğini (seçim yanlılığı) Golub verisi üzerinde göstermiştir. Düzeltilmiş dış 10-kat CV ile hata yaklaşık \%5 bulunmuştur. Varma ve Simon~\cite{varma2006}, hiperparametre seçimi de CV içinde yapıldığında yanlılığın sürdüğünü ve nested CV'nin gerçek hatayı neredeyse yansız tahmin ettiğini göstermiştir. Kapoor ve Narayanan~\cite{kapoor2023} ise veri sızıntısının günümüzde de 17 bilim alanında yaygın olduğunu raporlamıştır.

\subsection{Son Dönem Çalışmalar}
Son beş yılda metasezgisel ve hibrit öznitelik seçimi yöntemleri öne çıkmıştır~\cite{almazrua2022,elemam2022,osama2023,yaqoob2024}. Derin öğrenme tabanlı yaklaşımlar da önerilmiştir~\cite{alharbi2023,mallick2023}. Ancak tablo verisinde ağaç tabanlı modellerin derin modellerden çoğu zaman üstün olduğu bilinmektedir~\cite{grinsztajn2022}. Tablo verisi için transformer tabanlı bir temel model olan TabPFN~\cite{hollmann2025} küçük örneklemde güçlüdür ama yaklaşık 500 öznitelik sınırı nedeniyle gen ifadesi verisinde önce öznitelik seçimi gerektirir~\cite{zhou2025}. GenePT~\cite{chen2024} ise gen açıklamalarının LLM gömmelerini kullanan alternatif bir temsil sunmaktadır.

\subsection{Literatür Özet Tablosu}
Golub verisini kullanan çalışmaların özeti Tablo~\ref{tab:lit}'de verilmiştir.

\begin{table*}[t]
\caption{Golub verisini kullanan çalışmaların karşılaştırmalı özeti}
\label{tab:lit}
\centering
\footnotesize
\begin{tabularx}{\textwidth}{@{}l L l l l L L@{}}
\toprule
Yazar/Yıl & Ön işleme / öznitelik seçimi & Yöntem & Doğrulama & Seçim CV içinde mi? & En iyi sonuç & Not \\
\midrule
Golub 1999~\cite{golub1999} & SNR, 50 gen & Ağırlıklı oylama & LOOCV + 38/34 test & İçinde & 34 testin 29'unda kesin tahmin, tamamı doğru & 5 örnek belirsiz \\
Guyon 2002~\cite{guyon2002} & Standardizasyon; SVM-RFE, 2--16 gen & Doğrusal SVM & LOOCV + test & Dışında & 8--16 genle test hatası 0 & Seçim yanlılığı riski \\
Ambroise 2002~\cite{ambroise2002} & SVM-RFE, SFS & SVM, FLD & Dış 10-kat CV, .632+ & İçinde & $\sim$\%5 hata & Yanlılık düzeltmesi \\
Dudoit 2002~\cite{dudoit2002} & Kırpma 100--16.000, log$_{10}$; BW, 40 gen & DLDA, k-NN, CART vb. & 2:1 bölme $\times$200 & İçinde & Medyan 1 test hatası & Basit yöntemler yeterli \\
Díaz-Uriarte 2006~\cite{diaz2006} & Filtre yok; RF önem skoru, 2 gen & RF & .632+ bootstrap & İçinde & .632+ hata 0,075 & Gen listesi kararsız \\
Elemam 2022~\cite{elemam2022} & $\chi^2$, F, MI + SFS, 5 gen & SVM, KNN, DT, RF & 10-kat CV & Kısmen dışında & \%100 & Yalnızca doğruluk \\
Mallick 2023~\cite{mallick2023} & Seçim yok & 5 katmanlı DNN & 80/20 sabit bölme & Uygulanamaz & \%98,2--100 & $\sim$15 test örneği \\
Yaqoob 2024~\cite{yaqoob2024} & RDO, 4--14 gen & XGBoost & 3-kat CV, LOOCV & Dışında & \%99,87--100 & Örnek sayısında iç tutarsızlık \\
\bottomrule
\end{tabularx}
\end{table*}

\subsection{Literatürdeki Boşluk}
Golub verisinde yayımlanmış sonuçların çoğu \%95--100 doğruluk bildirmektedir. İncelenen sekiz çalışmanın hiçbiri nested CV kullanmamış, en az üçünde öznitelik seçimi kısmen ya da tamamen CV dışında yapılmıştır. Sınıflar dengesiz olduğu halde çoğu çalışma yalnızca doğruluk raporlamıştır. Sızıntıyı düzelten çalışmalarda~\cite{ambroise2002,diaz2006} hata \%5--7,5 düzeyine çıkmaktadır. Ayrıca hiçbir çalışma bulguları farklı bir platformda üretilmiş bağımsız bir veri setinde test etmemiştir; Golub'un test seti de aynı yongadan gelmektedir. Bu çalışma, sızıntısız nested CV, dengesizliğe duyarlı metrikler ve platformlar arası dış doğrulama ile bu boşluğu hedeflemektedir.

% =====================================================================
\section{Veri Setleri}
Ana veri seti olarak Golub ve ark.~\cite{golub1999} lösemi verisi, dış doğrulama için ise MILE çalışmasının Aşama~I verisi (GEO GSE13159)~\cite{haferlach2010} seçilmiştir.

\subsection{Golub Verisi}
Veri, Kaggle üzerinde CC0 (kamu malı) lisansıyla yayımlanan sürümden alınmıştır. Kemik iliği ve periferik kan örneklerinden oluşan 72 hastalık bu set, Affymetrix Hu6800 yongası (7.129 prob, 6.817 insan geni) ile ölçülmüştür. Referans bölünme 38 eğitim (27 ALL, 11 AML) ve 34 bağımsız test (20 ALL, 14 AML) örneğidir. Her hasta sütununun yanında Affymetrix'in P/A/M (Present/Absent/Marginal) tespit çağrısı yer alır ve hasta sütunları dosyada numara sırasında değildir.

\subsection{MILE Verisi (GSE13159)}
Microarray Innovations in LEukemia (MILE) çalışması, üç kıtadaki 11 merkezde yürütülmüş çok merkezli bir programdır~\cite{haferlach2010}. Aşama~I verisi 2.096 hastadan oluşur ve 18 tanı sınıfı içerir. Örnekler Ficoll ile saflaştırılmış mononükleer hücrelerden HG-U133 Plus 2.0 yongasıyla ölçülmüştür. Ön aşamada merkezler arası tekrarlanabilirliğin yüksek olduğu gösterilmiştir~\cite{kohlmann2008}. Bu çalışmada yalnızca ALL (C1--C8 alt sınıfları, $n=750$) ve AML alt sınıfları kullanılacaktır.

\subsection{Veri Setlerinin Karşılaştırması}
İki veri seti Tablo~\ref{tab:veri}'de yan yana özetlenmiştir.

\begin{table}[htbp]
\caption{Veri setlerinin karşılaştırması}
\label{tab:veri}
\centering
\footnotesize
\begin{tabularx}{\columnwidth}{@{}l L L@{}}
\toprule
Özellik & Golub (1999) & MILE Aşama I \\
\midrule
Erişim & Kaggle & GEO GSE13159 \\
Lisans & CC0 & GEO kamuya açık, atıf zorunlu \\
Merkez & Tek merkez & 3 kıtada 11 merkez \\
Doku & Kemik iliği, periferik kan & Kemik iliği, periferik kan (Ficoll) \\
Platform & Hu6800 (GPL80) & HG-U133 Plus 2.0 (GPL570) \\
$p$ & 7.129 prob & 54.613 prob seti \\
$n$ & 72 (38 / 34) & ALL+AML: \todo{koddan} \\
Sınıflar & ALL 47 / AML 25 & ALL 750 / AML \todo{koddan} \\
Ölçek & Ham MAS sinyali, negatif değerler var & Normalize: \todo{doğrulanacak} \\
\bottomrule
\end{tabularx}
\end{table}

\subsection{Seçim Gerekçesi ve Uyumluluk}
MILE dört nedenle seçilmiştir. Birincisi, aynı klinik problemi (ALL/AML) içerir ve ALL grubu Golub'daki gibi hem B hem T soy ALL örneklerini kapsar. İkincisi, aynı üreticinin farklı nesil yongasıyla ölçüldüğü için platformlar arası genellenebilirliği test etmeye uygundur. Üçüncüsü, çok merkezli tasarımı gerçekçi bir batch etkisi içerir. Dördüncüsü, büyük örneklemi Golub'un $p \gg n$ rejiminin tersini sunar.

Uyumlaştırma için iki platformun prob kimlikleri GPL80 ve GPL570 anotasyonlarıyla gen sembollerine eşlenecek ve ortak genler üzerinde çalışılacaktır. Ortak gen sayısı: \todo{koddan}. İki sete aynı dönüşüm (kırpma ve log$_2$) uygulanacaktır. Sınıf oranlarının farklı olması ve MILE'daki alt tip heterojenliği yorumlarda dikkate alınacaktır.

% =====================================================================
\section{Keşifsel Veri Analizi (EDA)}
\todo{Bu bölüm kod çıktıları geldikçe doldurulacaktır.}
\subsection{Veri Kalitesi}
\todo{Eksik/INF değerler, negatif ve uç değerler, P/A/M dağılımı, etiket eşleştirme doğrulaması.}
\subsection{Betimsel İstatistikler ve Sınıf Dağılımı}
\todo{Log dönüşümü öncesi ve sonrası ifade dağılımları.}
\subsection{Boyut İndirgeme}
\todo{PCA (açıklanan varyans), t-SNE / UMAP.}
\subsection{Ayırt Edici Genler}
\todo{FDR düzeltmeli sınıflar arası testler, en iyi 50 gen ısı haritası, hiyerarşik kümeleme.}
\subsection{Batch Analizi}
\todo{Golub eğitim--test karşılaştırması; Golub--MILE dağılım farkları.}
\subsection{Bulguların Özeti ve Modelleme İçin Çıkarımlar}
\todo{Numaralı bulgular: ölçekleme, öznitelik seçimi, dengesizlik.}

% =====================================================================
\section{Finale Kadar Plan}
Finalde üç yöntem, aynı sızıntısız pipeline içinde (imputasyon $\rightarrow$ kırpma ve log$_2$ $\rightarrow$ varyans filtresi $\rightarrow$ öznitelik seçimi $\rightarrow$ ölçekleme $\rightarrow$ model) ve aynı nested CV katlarında karşılaştırılacaktır.

\subsection{Seçilmesi Planlanan Yöntemler}
Planlanan yöntemler Tablo~\ref{tab:yontem}'de verilmiştir. Modülde verilen karar ağacı, taban çizgisi olarak ayrıca raporlanacaktır.

\begin{table}[htbp]
\caption{Planlanan yöntemler}
\label{tab:yontem}
\centering
\footnotesize
\begin{tabularx}{\columnwidth}{@{}l L l@{}}
\toprule
Yöntem & Gerekçe & Hiperparametreler \\
\midrule
Elastic-Net LR & Doğrusal, yorumlanabilir; gömülü öznitelik seçimi; $p \gg n$ için güçlü taban çizgisi & $C$, l1\_ratio \\
SVM (doğrusal, RBF) & Golub literatüründe en yaygın yöntem~\cite{furey2000,guyon2002}; doğrudan karşılaştırma & $C$, $\gamma$ \\
Rastgele Orman & Doğrusal olmayan, ölçeklemeye duyarsız; öznitelik önem skoru~\cite{diaz2006} & n\_estimators, max\_features \\
\bottomrule
\end{tabularx}
\end{table}

\subsection{Hedeflenen Bonuslar}
\begin{itemize}[leftmargin=*]
  \item \textbf{B1 -- Derin öğrenme:} dropout ve weight decay'li MLP; aynı dış CV katlarıyla, en az 5 seed.
  \item \textbf{B2 -- Transformer:} TabPFN v2~\cite{hollmann2025}; 500 öznitelik sınırı nedeniyle iç katlarda öznitelik seçimiyle birlikte~\cite{zhou2025}.
  \item \textbf{B4 -- Üçüncü veri seti (opsiyonel):} Armstrong ve ark.~\cite{armstrong2002} verisi, süre kalırsa.
\end{itemize}

\subsection{Takvim}
\begin{table}[htbp]
\caption{Finale kadar takvim}
\centering
\footnotesize
\begin{tabularx}{\columnwidth}{@{}l L@{}}
\toprule
Hafta & İş \\
\midrule
Ara rapor & EDA, ara rapor ve ara video teslimi \\
+1 & Pipeline, nested CV altyapısı, sabit 38/34 bölme deneyleri \\
+2 & Üç yöntemle ana deneyler, hiperparametre arama, istatistiksel test \\
+3 & MILE'a aynı protokolün uygulanması \\
+4 & Bonus deneyleri (B1, B2) \\
+5 & Yorumlanabilirlik, ablasyon, IEEE raporu ve final videosu \\
\bottomrule
\end{tabularx}
\end{table}

\subsection{Riskler}
\begin{table}[htbp]
\caption{Riskler ve önlemler}
\centering
\footnotesize
\begin{tabularx}{\columnwidth}{@{}L L@{}}
\toprule
Risk & Önlem \\
\midrule
Golub--MILE ortak gen sayısının düşük çıkması & Gen sembolü düzeyinde eşleme; ortak genlerle sınırlı analiz \\
Platformlar arası batch etkisi & Her veri setinde ayrı nested CV; örnek içi standardizasyon denemesi \\
Küçük örneklemde yüksek varyans & Tekrarlı CV ($\geq$5 seed), güven aralıkları, eşleştirilmiş istatistiksel test \\
Derin modellerin aşırı öğrenmesi & Erken durdurma iç katlarda; sonuçların dürüstçe raporlanması~\cite{grinsztajn2022} \\
MILE'ın büyük boyutu & Varyans filtresiyle boyut azaltma, paralel CV \\
\bottomrule
\end{tabularx}
\end{table}

% =====================================================================
\section*{Teşekkür ve GenAI Beyanı}
Bu raporun hazırlanmasında Anthropic Claude Opus 5.5 ve Google NotebookLM araçlarından yararlanılmıştır. Ayrıntılı kullanım beyanı ekte sunulmuştur. \todo{Ek GenAI beyanına göre güncellenecek.}

% =====================================================================
\begin{thebibliography}{99}
\bibitem{golub1999} T. R. Golub \textit{et al.}, ``Molecular classification of cancer: Class discovery and class prediction by gene expression monitoring,'' \textit{Science}, vol. 286, no. 5439, pp. 531--537, 1999, doi: 10.1126/science.286.5439.531.
\bibitem{furey2000} T. S. Furey \textit{et al.}, ``Support vector machine classification and validation of cancer tissue samples using microarray expression data,'' \textit{Bioinformatics}, vol. 16, no. 10, pp. 906--914, 2000, doi: 10.1093/bioinformatics/16.10.906.
\bibitem{guyon2002} I. Guyon, J. Weston, S. Barnhill, and V. Vapnik, ``Gene selection for cancer classification using support vector machines,'' \textit{Mach. Learn.}, vol. 46, pp. 389--422, 2002, doi: 10.1023/A:1012487302797.
\bibitem{dudoit2002} S. Dudoit, J. Fridlyand, and T. P. Speed, ``Comparison of discrimination methods for the classification of tumors using gene expression data,'' \textit{J. Amer. Stat. Assoc.}, vol. 97, no. 457, pp. 77--87, 2002, doi: 10.1198/016214502753479248.
\bibitem{ambroise2002} C. Ambroise and G. J. McLachlan, ``Selection bias in gene extraction on the basis of microarray gene-expression data,'' \textit{Proc. Natl. Acad. Sci. USA}, vol. 99, no. 10, pp. 6562--6566, 2002, doi: 10.1073/pnas.102102699.
\bibitem{saeys2007} Y. Saeys, I. Inza, and P. Larrañaga, ``A review of feature selection techniques in bioinformatics,'' \textit{Bioinformatics}, vol. 23, no. 19, pp. 2507--2517, 2007, doi: 10.1093/bioinformatics/btm344.
\bibitem{diaz2006} R. Díaz-Uriarte and S. Alvarez de Andrés, ``Gene selection and classification of microarray data using random forest,'' \textit{BMC Bioinformatics}, vol. 7, art. 3, 2006, doi: 10.1186/1471-2105-7-3.
\bibitem{varma2006} S. Varma and R. Simon, ``Bias in error estimation when using cross-validation for model selection,'' \textit{BMC Bioinformatics}, vol. 7, art. 91, 2006, doi: 10.1186/1471-2105-7-91.
\bibitem{kapoor2023} S. Kapoor and A. Narayanan, ``Leakage and the reproducibility crisis in machine-learning-based science,'' \textit{Patterns}, vol. 4, no. 9, art. 100804, 2023, doi: 10.1016/j.patter.2023.100804.
\bibitem{haferlach2010} T. Haferlach \textit{et al.}, ``Clinical utility of microarray-based gene expression profiling in the diagnosis and subclassification of leukemia: Report from the International Microarray Innovations in Leukemia Study Group,'' \textit{J. Clin. Oncol.}, vol. 28, no. 15, pp. 2529--2537, 2010, doi: 10.1200/JCO.2009.23.4732.
\bibitem{kohlmann2008} A. Kohlmann \textit{et al.}, ``An international standardization programme towards the application of gene expression profiling in routine leukaemia diagnostics: The Microarray Innovations in LEukemia study prephase,'' \textit{Br. J. Haematol.}, vol. 142, no. 5, pp. 802--807, 2008, doi: 10.1111/j.1365-2141.2008.07261.x.
\bibitem{armstrong2002} S. A. Armstrong \textit{et al.}, ``MLL translocations specify a distinct gene expression profile that distinguishes a unique leukemia,'' \textit{Nat. Genet.}, vol. 30, no. 1, pp. 41--47, 2002, doi: 10.1038/ng765.
\bibitem{almazrua2022} H. AlMazrua and H. Alshamlan, ``A comprehensive survey of recent hybrid feature selection methods in cancer microarray gene expression data,'' \textit{IEEE Access}, vol. 10, pp. 71427--71449, 2022, doi: 10.1109/ACCESS.2022.3185226.
\bibitem{elemam2022} T. Elemam and M. Elshrkawey, ``A highly discriminative hybrid feature selection algorithm for cancer diagnosis,'' \textit{Sci. World J.}, vol. 2022, art. 1056490, 2022, doi: 10.1155/2022/1056490.
\bibitem{grinsztajn2022} L. Grinsztajn, E. Oyallon, and G. Varoquaux, ``Why do tree-based models still outperform deep learning on typical tabular data?,'' in \textit{Proc. NeurIPS}, 2022, doi: 10.48550/arXiv.2207.08815.
\bibitem{osama2023} S. Osama, H. Shaban, and A. A. Ali, ``Gene reduction and machine learning algorithms for cancer classification based on microarray gene expression data: A comprehensive review,'' \textit{Expert Syst. Appl.}, vol. 213, art. 118946, 2023, doi: 10.1016/j.eswa.2022.118946.
\bibitem{alharbi2023} F. Alharbi and A. Vakanski, ``Machine learning methods for cancer classification using gene expression data: A review,'' \textit{Bioengineering}, vol. 10, no. 2, art. 173, 2023, doi: 10.3390/bioengineering10020173.
\bibitem{ilyas2023} M. Ilyas, K. M. Aamir, S. Manzoor, and M. Deriche, ``Linear programming based computational technique for leukemia classification using gene expression profile,'' \textit{PLOS ONE}, vol. 18, no. 10, art. e0292172, 2023, doi: 10.1371/journal.pone.0292172.
\bibitem{mallick2023} P. K. Mallick \textit{et al.}, ``Convergent learning--based model for leukemia classification from gene expression,'' \textit{Pers. Ubiquitous Comput.}, vol. 27, pp. 1103--1110, 2023, doi: 10.1007/s00779-020-01467-3.
\bibitem{yaqoob2024} A. Yaqoob \textit{et al.}, ``Optimizing cancer classification: A hybrid RDO-XGBoost approach for feature selection and predictive insights,'' \textit{Cancer Immunol. Immunother.}, vol. 73, no. 12, art. 261, 2024, doi: 10.1007/s00262-024-03843-x.
\bibitem{chen2024} Y. Chen and J. Zou, ``Simple and effective embedding model for single-cell biology built from ChatGPT,'' \textit{Nat. Biomed. Eng.}, 2024, doi: 10.1038/s41551-024-01284-6.
\bibitem{hollmann2025} N. Hollmann \textit{et al.}, ``Accurate predictions on small data with a tabular foundation model,'' \textit{Nature}, vol. 637, no. 8045, pp. 319--326, 2025, doi: 10.1038/s41586-024-08328-6.
\bibitem{zhou2025} \todo{Yazarlar PDF'ten doğrulanacak}, ``The limitations of TabPFN for high-dimensional RNA-seq analysis,'' \textit{bioRxiv} (ön baskı), 2025, doi: 10.1101/2025.08.15.670537.
\end{thebibliography}

\end{document}
